\documentclass[11pt]{article}
\usepackage[margin=1in]{geometry}
\usepackage{amsmath,amssymb,amsthm}
\usepackage{booktabs}
\usepackage[numbers,sort&compress]{natbib}
\usepackage[expansion=false]{microtype}
\usepackage{xcolor}
\usepackage[colorlinks=true,linkcolor=blue!50!black,citecolor=blue!50!black,urlcolor=blue!50!black]{hyperref}
\hypersetup{pdftitle={New lower bounds for trifferent codes of lengths 10, 11, 12, 16 and 17},pdfauthor={Tyler Duic}}

\newtheorem{theorem}{Theorem}
\newtheorem{lemma}[theorem]{Lemma}
\theoremstyle{remark}
\newtheorem{remark}[theorem]{Remark}

\newcommand{\F}{\mathbb{F}}
\newcommand{\w}[1]{\texttt{#1}}

\title{New lower bounds for trifferent codes\\ of lengths 10, 11, 12, 16 and 17}
\author{Tyler Duic\\[2pt]\small Independent researcher, New York, NY\\\small ORCID \href{https://orcid.org/0009-0005-4675-3045}{0009-0005-4675-3045}\\\small\url{https://github.com/tduic}}
\date{October 2026}

\begin{document}
\maketitle

\begin{abstract}
A code $C\subseteq\{0,1,2\}^n$ is trifferent if any three distinct codewords have a coordinate in which they are pairwise distinct; equivalently, $C$ is a perfect hash family $\mathrm{PHF}(n;|C|,3,3)$.
Writing $T(n)$ for the largest size of a trifferent code of length $n$, I give explicit codes showing
$T(10)\ge 30$, $T(11)\ge 36$, $T(12)\ge 48$, $T(16)\ge 108$ and $T(17)\ge 109$.
Each improves on every bound I found in the literature, including the closure of the known values under the elementary monotonicity and shortening rules.
The codes have short algebraic descriptions: unions of orbits of a translation and a cyclic shift for $n\le 12$, and four cosets of a linear $[16,3]_3$ code for $n=16$.
Each bound is verified by the Lean~4 proof assistant against a direct definition of trifference.
\end{abstract}

\section{Introduction}

A code $C\subseteq\{0,1,2\}^n$ is \emph{trifferent} if for any three distinct $x,y,z\in C$ there is a coordinate $i$ with $\{x_i,y_i,z_i\}=\{0,1,2\}$.
Let $T(n)$ denote the maximum size of a trifferent code of length $n$.
Determining the growth rate of $T(n)$ is a long-standing problem in information theory and combinatorics; the classical bounds are due to Fredman and Koml\'os~\cite{fredman1984} and K\"orner and Marton~\cite{korner1988}.
Equivalently, a trifferent code of length $n$ and size $k$ is a perfect hash family $\mathrm{PHF}(n;k,3,3)$: the $n$ coordinates are functions from the $k$ codewords to a $3$-set, and every $3$-subset of codewords is separated by one of them~\cite{walker2007}.

Exact values are known for small lengths: $T(5)=10$ and $T(6)=13$~\cite{dellafiore2022}, and $T(7)=16$, $T(8)=20$, $T(9)=27$~\cite{kurz2024}.
For larger $n$, lower bounds come from tables of perfect hash families~\cite{walker2007} and from linear codes: a linear code over $\F_3$ is trifferent if and only if it is minimal, and minimal $[14,4]_3$ and $[19,5]_3$ codes give $T(14)\ge 81$ and $T(19)\ge 243$~\cite{bishnoi2023,kurz2024}.

This note gives five improved lower bounds for small lengths, each witnessed by an explicit code and checked by a proof assistant.

\begin{theorem}\label{thm:main}
$T(10)\ge 30$, $T(11)\ge 36$, $T(12)\ge 48$, $T(16)\ge 108$ and $T(17)\ge 109$.
\end{theorem}

\section{The baseline}

Two elementary rules propagate bounds between lengths.
First, $T(n)\ge T(n-1)+1$~\cite[Lemma~1 and Corollary~2]{kurz2024}: given a trifferent code $C$ of length $n-1$ and $c\in C$, the code $\{(c,0),(c,1)\}\cup\{(c',2):c'\in C\setminus\{c\}\}$ is trifferent of length $n$.
Second, \emph{shortening}: $T(n-1)\ge\lceil 2T(n)/3\rceil$, equivalently $T(n)\le\lfloor 3T(n-1)/2\rfloor$ (cf.~\cite{kurz2024}).
Indeed, if $C$ has length $n$, some symbol $a$ occurs in a given coordinate in at most $|C|/3$ codewords; deleting those codewords and then that coordinate leaves a trifferent code of length $n-1$ with at least $\lceil 2|C|/3\rceil$ words. Three remaining words take at most two values in the deleted coordinate, so they are separated by another coordinate; for the same reason (when at least three words remain) no two remaining words coincide after the deletion.

As a baseline I take the best published values---the exact values above, the linear-code bounds $T(14)\ge 81$ and $T(19)\ge 243$, and the best values in Table~6 of Walker and Colbourn~\cite{walker2007}, which are $29,30,36,37,48,57,81,82,83$ for $n=10,\dots,18$---and close them under both rules.
Table~\ref{tab:bounds} compares the result with Theorem~\ref{thm:main}.

\begin{table}[ht]
\centering
\begin{tabular}{rrlrr}
\toprule
$n$ & baseline & source of baseline & new bound & gain \\
\midrule
10 & 29 & Walker--Colbourn~\cite{walker2007} & 30 & $+1$\\
11 & 30 & Walker--Colbourn~\cite{walker2007} & 36 & $+6$\\
12 & 36 & Walker--Colbourn~\cite{walker2007} & 48 & $+12$\\
16 & 83 & $T(14)\ge 81$ and two steps of $T(n)\ge T(n-1)+1$ & 108 & $+25$\\
17 & 108 & shortening $T(18)\ge 162$, itself shortened from $T(19)\ge 243$ & 109 & $+1$\\
\bottomrule
\end{tabular}
\caption{Best bounds I found in the literature (the closure described in the text) and the new bounds.}
\label{tab:bounds}
\end{table}

\begin{remark}
I also found codes of size $54$ at $n=13$ and $162$ at $n=18$, but these sizes equal the shortenings of the minimal $[14,4]_3$ and $[19,5]_3$ codes, so they are not new bounds.
\end{remark}

\section{The codes}

Throughout, codewords are written as strings over $\{0,1,2\}$, and for $s\mid n$ let $G_s$ be the group of maps on $\{0,1,2\}^n$ generated by $x\mapsto x+(1,\dots,1)\pmod 3$ and the cyclic shift of coordinates by $s$ positions.

\paragraph{Length 10.} The code is the union of the $G_2$-orbits of
\w{0001101220} and \w{0012020120}, each of size $15$, giving $30$ words.

\paragraph{Length 11.} The code is the union of the $G_1$-orbits of
\w{00000000000} (size $3$) and \w{00011101201} (size $33$), giving $36$ words.

\paragraph{Length 12.} The code is the union of the $G_2$-orbits of
\w{000111222111}, \w{000121012221} (each of size $18$), \w{000122201112} and \w{010201020102} (each of size $6$), giving $48$ words.

\paragraph{Length 16.} Let $D\subseteq\F_3^{16}$ be the linear code spanned by
\[
\w{0001221202201001},\qquad \w{0110120102221111},\qquad \w{1010100010111101}.
\]
It has dimension $3$ and nonzero weights $9$ and $12$; being trifferent and linear, it is a minimal code~\cite{bishnoi2023}.
The code is the union of the four cosets $r+D$ with
\[
r\in\{\w{0000000000000000},\ \w{0010011222101021},\ \w{0010122101011012},\ \w{0020221012102011}\},
\]
giving $108$ words.

\paragraph{Length 17.} Apply the construction of~\cite[Lemma~1]{kurz2024} to the length-16 code with $c=0$: the code is $\{(0,0),(0,1)\}\cup\{(c',2): c'\in C_{16}\setminus\{0\}\}$, where $C_{16}$ is the length-16 code, with $109$ words.

\medskip
The full codes are listed in machine-readable form in the accompanying repository, together with a script that rebuilds them from the descriptions above.

\section{Verification}

Checking trifference directly is a finite computation over $\binom{|C|}{3}$ triples---about $2\times 10^5$ for the length-16 code---but I wanted a check that does not rely on trusting a particular program.
I therefore verified each bound in the Lean~4 proof assistant~\cite{demoura2021lean4} using Mathlib~\cite{mathlib2020}.

\begin{itemize}
\item \textbf{Definition.} Trifference is defined directly on finite sets of functions $\{1,\dots,n\}\to\{0,1,2\}$: for all pairwise distinct $x,y,z\in C$ there is $i$ with $x_i,y_i,z_i$ pairwise distinct.
\item \textbf{Statement.} Each bound in Theorem~\ref{thm:main} is proved in the form $\exists\, C\subseteq\{0,1,2\}^n,\ C$ trifferent $\wedge\ k\le |C|$.
\item \textbf{Method.} A bitmask checker on lists of codewords is proved sound once, in general; it is then evaluated on each explicit code by the Lean kernel (using \texttt{decide} with kernel reduction, in chunks to bound memory). No compiled or native evaluation is trusted, and the only axioms used are the standard ones of Lean's logic (\texttt{propext}, \texttt{Classical.choice}, \texttt{Quot.sound}).
\item \textbf{Independent statement check.} A separately written definition of trifference is proved equivalent to the one used, each bound is restated with it, and a three-word code that is not trifferent is shown to be rejected, so the definition is not vacuous.
\end{itemize}

Independently of Lean, every code was re-checked triple by triple with separate code, and the Lean certificates were decoded and compared word for word with the published codes.

\section{Remarks}

The gains at $n=10$ and $n=17$ are small; the latter is a direct consequence of the length-16 code.
The largest gains are at $n=12$ and $n=16$, where structured codes---orbit unions and cosets of a minimal linear code---exceed both the values in the perfect hash family tables and the linear-code bounds propagated by the elementary rules.
I compared against the trifferent-code literature cited above and the perfect hash family tables of~\cite{walker2007}. I was unable to consult later online tables of perfect hash families and would be grateful to learn of any entry there that matches or exceeds these values.

\paragraph{Data and code availability.}
The codes, the Lean~4 proofs and the checking scripts are at \url{https://github.com/tduic/research-results/tree/main/trifferent-codes}.

\paragraph{Acknowledgments.}
I carried out the search, the Lean formalization and the independent checks with substantial assistance from AI models (Anthropic's Claude), working under my direction. I checked every claim and reference in this note and take full responsibility for its content.

\bibliographystyle{unsrtnat}
\bibliography{references}

\end{document}
